\begin{abstract}
Local coordination and linker chemistry can organize differences between partial-charge definitions, but element offsets and nonlinear relations between charge scales can mimic an environment effect. We quantify the chemical contribution for DDEC6, CM5, and Bader charges in 9,171 QMOF structures containing 1,420,532 atoms. Matched predictors add first-shell chemistry and two-hop path composition to a flexible control supplied with the central element and source charge. Local information improves held-out all-atom and nonmetal predictions for all three charge pairs, including a split that blocks identical source MOFkeys or reduced formulas. Under this split, DDEC6--Bader structure-balanced mean absolute error decreases from 0.08454 to 0.03885~\charge. A paired chemical atlas shows that directional shifts and greater absolute departures are distinct: O-attached carbon has a DDEC6--Bader signed shift without an established increase in absolute departure, whereas N-attached carbon has a CM5--Bader magnitude increase with a near-zero signed contrast. Central-metal behavior is mixed, including worse DDEC6--Bader prediction after adding structure to the global model. Finite connectivity tests retain the tested nonmetal contrasts while exposing shared-proton ambiguity. The atlas thus identifies conditional chemical structure in charge-definition disagreement, with an explicit distinction between empirical departures and physical charge accuracy.
\end{abstract}
\noindent\textbf{Keywords:} metal--organic frameworks; partial atomic charges; local chemical environments; QMOF; conditional prediction
